\section{Experimental Design}

\paragraph{Positive-control gate.}
Before any novel pilot, the pipeline must reproduce the frozen two-arm comparison from Drayson et al.\ \cite{drayson2025machine}: the official fully synthetic recursive condition and the official human-mixed condition. The upstream repository commit, dependency environment, GPT-2 model and tokenizer revisions, dataset revision, decoding configuration, horizon, endpoint, expected ordering, tolerance, and rerun rule must be recorded before execution. Both arms use one unchanged comparison protocol. A pass requires completion of every planned generation, recovery of the predeclared direction, satisfaction of the frozen endpoint criterion, and a clean rerun with the same scientific conclusion. If this gate fails, commands, logs, configs, manifests, hashes, deviations, and partial artifacts are preserved as a failure package; the novel pilot remains blocked.

\paragraph{Domain, model, horizon, and budgets.}
The August study is intentionally a screening-scale experiment: one legally usable text domain, one 124M--160M language model, one primary nonzero lifetime human-token budget, and a recursive horizon supported by the compute benchmark. The current planning value is ten generations, but it is not final until the preregistration records the exact model, tokenizer, initialization rule, context length, optimizer, update count, decoding parameters, synthetic-pool size, $G$, human budget $B$, and total optimizer budget $T$. All policies in the primary contrast use the same values. \textbf{PENDING\_DATA\_FREEZE:} the final licensed domain and manifest hashes await Neil's signed audit. \textbf{PENDING\_METHOD\_FREEZE:} horizon, budget, seeds, and policy configs await Aarav's freeze. These fields cannot be inferred from favorable treatment results.

\paragraph{Five disjoint partitions.}
The study uses exactly five manifest-backed partitions: base human training data, human rescue candidates, generation prompts, validation data, and untouched final human test data. Each example carries a stable identifier, normalized content hash, source dataset and revision, license/provenance metadata, and preprocessing version. Exact duplicates are forbidden across partitions; the frozen near-duplicate procedure reports and resolves prohibited overlap before a run. Final-test text is unavailable to prompting, mode construction, policy scoring, threshold selection, early stopping, hyperparameter selection, and screening. Full licensed corpora and large generations remain outside Git; manifests and artifact references store immutable locations, SHA-256 hashes, sizes, and producing commits.

\paragraph{Paired recursive chains.}
For every seed, the compared policies begin from the same validated base checkpoint and use the same manifests, prompt order, evaluation data, and frozen generation/evaluation settings. Random streams are deterministically derived for initialization, data sampling, generation, dropout, policy selection, and evaluation. The seed order and replacement rule are preregistered. A complete chain, not a generation, is the experimental unit. Interrupted chains may resume only through the checkpoint contract and must satisfy the frozen equivalence rule. Hardware failures, missing artifacts, and corrupt hashes are implementation failures; numerical divergence is a scientific outcome unless an independently verifiable defect caused it. All failures remain in the failure log.

\paragraph{Primary and secondary outcomes.}
Held-out human negative log-likelihood is computed from target-token log probabilities on the final-test manifest, excluding padding and recording the evaluated token count. Generation-wise NLL regret is aggregated over the fixed horizon using preregistered weights to form one chain-level area-under-the-curve outcome. The second primary outcome is one tail-retention metric frozen only after deterministic, finite-input, directionality, and independence checks. It may not equal the policy's selection score. Secondary outcomes may include per-mode regret, semantic coverage, generated-text diversity, overlap/memorization diagnostics, allocation trajectories, and compute efficiency; they cannot replace a failed primary outcome.

\paragraph{Analysis and interpretation.}
The planned estimand is the paired chain-level difference between the joint policy and the strongest eligible non-joint policy identified by the frozen screening rule. Uncertainty is calculated across independent chains; treating generations as independent replicates is prohibited. The primary budget, horizon, contrast, meaningful-effect threshold, multiplicity treatment, missing-run rule, and seed count are fixed before confirmation outcomes. A favorable conclusion requires valid manifests, exact $B$ and $T$ matching, no leakage, the preregistered effect criterion, and uncertainty excluding the equivalent or harmful region. A tie, harmful result, or wide interval is reported directly. The monitoring-omission condition is a prespecified boundary test, not a license to retune the main monitor.

\paragraph{Compute and artifact discipline.}
Pilot estimates record expected accelerator-hours, wall time, peak memory when available, storage, and artifact destinations before launch; actual use is appended without erasing history. Every completed chain preserves the repository commit and working-tree state, environment, config hashes, manifest hashes, checkpoint and generated-data hashes, allocation history, metrics, stdout/stderr, status, and failure reason. Tables and figures are generated from validated immutable chain results and receive content hashes. \textbf{RESULT\_PENDING:} no novel value or interpretation belongs in this section; any unresolved protocol item remains an explicit pending marker until frozen by its owner.
